\documentclass[12pt,a4paper]{article}
\usepackage{amsmath,amssymb,amsthm}
\usepackage{hyperref}
\usepackage{geometry}
\geometry{margin=2.5cm}
\usepackage{enumitem}
\usepackage{booktabs}

\newtheorem{theorem}{Theorem}[section]
\newtheorem{lemma}[theorem]{Lemma}
\newtheorem{corollary}[theorem]{Corollary}
\newtheorem{proposition}[theorem]{Proposition}
\theoremstyle{definition}
\newtheorem{definition}[theorem]{Definition}
\theoremstyle{remark}
\newtheorem{remark}[theorem]{Remark}

\title{Hawking Radiation Temperature from Recognition Science:\\
Closed-Form Derivation in $\varphi$-Ladder Units}

\author{Jonathan Washburn\\
Recognition Science Research Institute\\
Austin, Texas\\
\texttt{washburn.jonathan@gmail.com}}

\date{March 2026}

\begin{document}
\maketitle

\begin{abstract}
We derive the Hawking radiation temperature from the Recognition
Science (RS) forcing chain.  The standard formula $T_H = \hbar c^3/
(8\pi G M k_B)$ is expressed entirely in terms of the golden ratio
$\varphi$ and the black hole mass $M$.  In RS-native units
($\hbar = \varphi^{-5}$, $G = \varphi^5$, $c = 1$,
$k_R = \ln\varphi$):
\begin{equation*}
  T_H = \frac{1}{8\pi\,\varphi^{10}\,M\,\ln\varphi}.
\end{equation*}
The Einstein coupling constant $\kappa = 8\pi G/c^4 = 8\pi\varphi^5$
enters through the surface gravity $\kappa_{\mathrm{surf}} =
c^4/(4GM)$.  We prove $T_H > 0$ for all $M > 0$, derive the
inverse-mass scaling $T_H \propto M^{-1}$, and compute the
temperature for representative black hole masses.
\end{abstract}

%% ============================================================
\section{Introduction}
%% ============================================================

Hawking's 1974 discovery~\cite{Hawking1974} that black holes radiate
at temperature $T_H \propto M^{-1}$ revealed a deep connection between
gravity, quantum mechanics, and thermodynamics.  The standard formula:
\begin{equation}
  T_H = \frac{\hbar c^3}{8\pi G M k_B},
  \label{eq:TH-SI}
\end{equation}
contains four fundamental constants ($\hbar$, $c$, $G$, $k_B$) and
one parameter ($M$).  In RS~\cite{RS_Axioms}, all four constants are
derived from the forcing chain, making $T_H$ a zero-parameter function
of mass alone.

%% ============================================================
\section{Recognition Science Framework}
\label{sec:framework}
%% ============================================================

Recognition Science derives all physics from a single primitive,
the \emph{recognition cost functional} $J(x)$, uniquely determined
by three axioms.

\begin{definition}[RS Cost Functional]
For $x > 0$: $J(x) = \tfrac{1}{2}(x + x^{-1}) - 1$.
\end{definition}

\noindent\textbf{Axiom A1 (Normalization):} $J(1) = 0$.\\
\textbf{Axiom A2 (Recognition Composition Law, RCL):}
$J(xy) + J(x/y) = 2J(x)J(y) + 2J(x) + 2J(y)$ for all $x, y > 0$.\\
\textbf{Axiom A3 (Calibration):} $J''_{\log}(0) = 1$, where
$J_{\log}(t) := J(e^t)$.

\begin{theorem}[Cost Uniqueness]
\label{thm:unique}
The continuous solution to \textup{(A1)--(A3)} is unique:
$J(x) = \tfrac{1}{2}(x + x^{-1}) - 1$.
\end{theorem}

\begin{proof}[Proof sketch]
Setting $f(t) = J_{\log}(t) + 1$ and rewriting \textup{(A2)} gives the
d'Alembert equation $f(s+t) + f(s-t) = 2f(s)f(t)$.  By the Acz\'el
classification theorem~\cite{Aczel1966}, the unique continuous solution
with $f(0) = 1$ and $f''(0) = 1$ is $f(t) = \cosh t$.
\end{proof}

\noindent The RS forcing chain yields the quantum-gravitational duality
$G \cdot \hbar = 1$, giving $\hbar = \varphi^{-5}$, $G = \varphi^5$,
$c = 1$, and $k_R = \ln\varphi$ in RS-native units.

%% ============================================================
\section{RS Constants}
\label{sec:constants}
%% ============================================================

\begin{definition}[RS Fundamental Constants]
From the RS forcing chain and the quantum-gravitational
duality $G \cdot \hbar = 1$~\cite{RS_Axioms}:
\begin{align}
  \hbar &= \varphi^{-5}\,\tau_0, \qquad
  c = \frac{\ell_0}{\tau_0}, \qquad
  G = \varphi^5\,\frac{\ell_0^2}{\tau_0}, \qquad
  k_R = \ln\varphi.
\end{align}
In RS-native units ($\ell_0 = \tau_0 = 1$):
$\hbar = \varphi^{-5}$, $c = 1$, $G = \varphi^5$
(so that $G\cdot\hbar = \varphi^5 \cdot \varphi^{-5} = 1$).
\end{definition}

\begin{definition}[Einstein Coupling]
$\kappa_{\mathrm{rs}} = 8\pi G/c^4 = 8\pi\varphi^5 \approx 278.6$.
\end{definition}

\begin{remark}[$\kappa_{\mathrm{rs}}$ Numerical Value]
Since $1.618 < \varphi < 1.619$, we have $1.618^5 \approx 11.09$
and $8\pi\varphi^5 \approx 8 \times 3.14159 \times 11.09 \approx 278.6$.
The Einstein coupling $\kappa_{\mathrm{rs}} = 8\pi\varphi^5$ lies
in the interval $(275, 282)$ at precision $\varphi \in (1.618, 1.619)$.
\end{remark}

%% ============================================================
\section{Derivation}
\label{sec:derivation}
%% ============================================================

\begin{theorem}[Hawking Temperature in RS-Native Units]
\label{thm:TH}
\begin{equation}
  T_H = \frac{1}{8\pi\,\varphi^{10}\,M\,\ln\varphi}.
\end{equation}
\end{theorem}

\begin{proof}
Substituting the RS constants ($\hbar = \varphi^{-5}$, $G = \varphi^5$,
$c = 1$, $k_R = \ln\varphi$) into~\eqref{eq:TH-SI}:
\begin{equation}
  T_H = \frac{\hbar\,c^3}{8\pi\,G\,M\,k_R}
  = \frac{\varphi^{-5} \cdot 1}{8\pi \cdot \varphi^5 \cdot M
  \cdot \ln\varphi}
  = \frac{1}{8\pi\,\varphi^{10}\,M\,\ln\varphi}.
\end{equation}
\end{proof}

\begin{theorem}[Positivity]
\label{thm:positive}
$T_H > 0$ for all $M > 0$.
\end{theorem}

\begin{proof}
$\varphi^{10} > 0$, $M > 0$, $\ln\varphi > 0$ (since $\varphi > 1$).
The denominator is the product of positive quantities, so
$T_H = 1/(8\varphi^{10} M \ln\varphi) > 0$.
\end{proof}

\begin{theorem}[Inverse-Mass Scaling]
$T_H \propto M^{-1}$: smaller black holes are hotter.
\end{theorem}

\begin{proof}
$T_H = C/M$ where $C = 1/(8\pi\varphi^{10}\ln\varphi) > 0$ is
independent of $M$.
\end{proof}

%% ============================================================
\section{Numerical Estimates}
\label{sec:numerical}
%% ============================================================

\begin{center}
\renewcommand{\arraystretch}{1.3}
\begin{tabular}{@{}lll@{}}
\toprule
\textbf{Black Hole} & \textbf{Mass} & \textbf{$T_H$} \\
\midrule
Solar mass & $M_\odot = 2 \times 10^{30}$~kg &
$6.2 \times 10^{-8}$~K \\
Primordial ($10^{12}$~kg) & $10^{12}$~kg &
$1.2 \times 10^{11}$~K \\
Planck mass & $m_P = 2.2 \times 10^{-8}$~kg &
$5.6 \times 10^{30}$~K \\
\bottomrule
\end{tabular}
\end{center}

\begin{remark}
For solar-mass black holes, $T_H$ is far below the CMB temperature
($2.725$~K), so the black hole absorbs more than it radiates and
grows.  Only black holes with $M \lesssim 10^{22}$~kg have $T_H$
above the CMB and can evaporate in the current epoch.
\end{remark}

%% ============================================================
\section{The Surface Gravity Connection}
\label{sec:surface-gravity}
%% ============================================================

\begin{definition}[Surface Gravity]
For a Schwarzschild black hole:
$\kappa_{\mathrm{surf}} = c^4/(4GM)$.
\end{definition}

\begin{proposition}[Temperature--Surface Gravity Relation]
\begin{equation}
  T_H = \frac{\hbar\,\kappa_{\mathrm{surf}}}{2\pi\,c\,k_R}
  = \frac{\kappa_{\mathrm{surf}}}{2\pi} \cdot
  \frac{\varphi^{-5}}{\ln\varphi}.
\end{equation}
\end{proposition}

\begin{proof}
Substituting $\kappa_{\mathrm{surf}} = c^4/(4GM) = 1/(4\varphi^5 M)$
(in RS-native units $c = 1$, $G = \varphi^5$):
\[
  \frac{\hbar\,\kappa_{\mathrm{surf}}}{2\pi\,c\,k_R}
  = \frac{\varphi^{-5} \cdot \frac{1}{4\varphi^5 M}}{2\pi \cdot \ln\varphi}
  = \frac{1}{8\pi\,\varphi^{10}\,M\,\ln\varphi} = T_H. \qedhere
\]
\end{proof}

\begin{remark}
This is the Unruh effect~\cite{Unruh1976} applied to the horizon:
an observer at rest near the horizon experiences proper acceleration
$\kappa_{\mathrm{surf}}$ and perceives a thermal bath at temperature
$T_H$.  In RS, the Unruh characteristic frequency is
$F_U = \kappa_{\mathrm{surf}}/(2\pi c) = c^3/(8\pi GM)$.
The Hawking luminosity (Stefan--Boltzmann applied to the horizon) is
\begin{equation}
  P_H = \frac{\hbar\,c^6}{15360\,\pi\,G^2\,M^2}
  = \frac{c^4\,F_U^2}{240\,\pi^2\,G\,M},
\end{equation}
where the numerical coefficient $1/15360\pi$ accounts for the
greybody factor for a single massless boson species.  Substituting
RS constants ($\hbar = \varphi^{-5}$, $c = 1$, $G = \varphi^5$):
\[
P_H = \frac{\varphi^{-5} \cdot 1}{15360\pi \cdot \varphi^{10} \cdot M^2}
= \frac{1}{15360\pi\,\varphi^{15}\,M^2}.
\]
(Here $G^2 = \varphi^{10}$ and $\hbar/G^2 = \varphi^{-5}/\varphi^{10}
= \varphi^{-15}$.)
\end{remark}

%% ============================================================
\section{Predictions}
\label{sec:predictions}
%% ============================================================

\begin{enumerate}
  \item \textbf{Exact $M^{-1}$ scaling}: $T_H \propto M^{-1}$ with no
  corrections from new physics.  Deviation would indicate additional
  degrees of freedom near the horizon.

  \item \textbf{No greybody factors in the RS limit}: The pure
  $\varphi$-ladder spectrum is thermal.  Greybody factors arise from
  the angular-momentum barrier of the external geometry, not from
  the emission process itself.

  \item \textbf{Temperature floor}: $T_H$ cannot exceed
  $T_{\max} = 1/(8\varphi^{10} M_{\min} \ln\varphi)$ where
  $M_{\min}$ is the Planck-mass remnant.
\end{enumerate}

%% ============================================================
\section{Conclusion}
\label{sec:conclusion}
%% ============================================================

The Hawking temperature is derived entirely from RS constants
($\hbar = \varphi^{-5}$, $G = \varphi^5$, $c = 1$, $k_R = \ln\varphi$):
$T_H = 1/(8\pi\varphi^{10}\,M\,\ln\varphi)$ in RS-native units.
No free parameters appear.  The temperature is provably
positive, scales as $M^{-1}$, and reproduces the standard
Hawking formula when expressed in SI units.  The Einstein coupling $\kappa_{\mathrm{rs}} = 8\pi\varphi^5$
enters through the surface gravity.

%% ============================================================
\begin{thebibliography}{99}

\bibitem{Aczel1966}
J.~Acz\'el, \emph{Lectures on Functional Equations and Their Applications},
Academic Press, New York (1966).

\bibitem{RS_Axioms}
J.~Washburn, ``The Algebra of Reality: A Recognition Science Derivation
of Physical Law,'' \emph{Axioms} \textbf{15}(2), 90 (2026).
\texttt{doi:10.3390/axioms15020090}

\bibitem{Hawking1974}
S.~W.~Hawking, ``Black Hole Explosions?'' Nature \textbf{248}, 30
(1974).

\bibitem{Unruh1976}
W.~G.~Unruh, ``Notes on Black-Hole Evaporation,'' Phys.\ Rev.\ D
\textbf{14}, 870 (1976).

\end{thebibliography}

\end{document}
